\section{Period recognition from the reconstructed patch}
\label{sec:periods}
This section retains the protected-patch argument for repeated shapes.
It acts on the acquired geometry, not on an assumed list of integer
copy labels. We include the proof to specify exactly which discrete
information the adaptive experiment recovers.

Put $B=1+L_0$ and $Q_a=[-a,a]^2$. The aperture used above is chosen
large enough to recover complete components with centers in $Q_{10B}$,
together with their protective collars. For two components define
\[
 d_*(C,D)=\|z_C-z_D\|_\infty+\|p_C-p_D\|_\infty,
\]
where the centered support functions are compared in the laboratory
directions, without an independent rotation. Define the patch defect
\begin{equation}\label{eq:defect}
 \mathcal D(w)=\max_{C:z_C\in Q_B}\ \max_{\epsilon\in\{-1,1\}}
                     \inf_D d_*(C+\epsilon w,D).
\end{equation}
The infimum ranges over physical components. Distant components cannot
minimize it. The known nonperiod margin is the requirement
\begin{equation}\label{eq:patch-margin}
 \begin{gathered}
 w=z_D-z_C,\quad z_C\in Q_{3B},\ z_D\in Q_{5B},\quad w\notin\Pi
       \quad\Longrightarrow\quad\mathcal D(w)\ge\eta>0.
 \end{gathered}
\end{equation}
This is a margin for whole-patch translations, not a separation between
all individual shapes. Every fixed table in $\mathcal B$ has some
positive such margin, since there are finitely many candidate pairs.
The margin is not claimed to be uniformly inferable from finite bits.

\begin{lemma}[Exact patch recognition]\label{lem:recognition}
One has $\mathcal D(w)=0$ if and only if $w\in\Pi$. For
$\|w\|_\infty\le8B$, this zero test uses only components centered in
$Q_B$ and $Q_{9B}$. The group $\Pi$ is a lattice with
\begin{equation}\label{eq:index-bound}
 [\Pi:L\Z^2]\le r_0,\qquad\covol\Pi\ge V_*=V_0/r_0.
\end{equation}
Choose any component centered in $Q_{2B}$ as a root. Its center
differences to components centered in $Q_{4B}$ that pass the zero
test generate $\Pi$ over $\Z$.
\end{lemma}
\begin{proof}
Each orbit modulo $L\Z^2$ has a representative centered in
$L[-1/2,1/2]^2\subset Q_B$. A zero defect gives both translates of
every such representative as physical components. Translating by
$L\Z^2$ gives $\mathcal O+w\subset\mathcal O$ and
$\mathcal O-w\subset\mathcal O$; together these imply equality.
The converse is immediate. The centers of a matching pair in the
bounded test are in $Q_{9B}$.

Map a coset of $\Pi/L\Z^2$ to the component orbit of the translated
root. This is injective: a nonempty compact body is not invariant
under a nonzero translation. Thus the quotient has at most $r$ elements.
A finite extension of a full-rank lattice is a lattice, giving
\eqref{eq:index-bound}. The two columns of $L$, and representatives
of all these period cosets, can be chosen with norm less than $B$
in the maximum norm. Translating the root by these vectors gives
targets in $Q_{3B}$, strictly within the cutoff $Q_{4B}$. They are
included in the accepted list and generate $\Pi$. Conversely each
accepted vector is a true period.
\end{proof}

\begin{lemma}[Stable discrete decisions]\label{lem:locking}
Suppose the protected complete components have matched Hausdorff
error at most $e$. On $\mathcal B_\eta$, for sufficiently small $e$
depending on the full prior, finite tests recover the true primitive
orbit count and exact rational relations of a true accepted generating
list in a true independent-pair basis. They recover a matched period
basis and primitive free area to error $O(e)$.
\end{lemma}
\begin{proof}
Hausdorff error $e$ gives Steiner-center error at most $2e$ and
centered-support error at most $3e$. Use approximate source centers
in $Q_{2B}$, a root there and root targets in $Q_{4B}$. For a measured
difference $\widehat w$, its corresponding true difference has error
at most $4e$. Test both signs on all selected sources, minimizing over
targets with approximate centers in $Q_{10B}$. The center contribution
to each comparison changes by at most $8e$ and the support contribution
by at most $6e$. True periods therefore have measured defect at most
$14e$. For a nonperiod, its true endpoints lie in $Q_{3B},Q_{5B}$,
so \eqref{eq:patch-margin} supplies a witness of defect at least $\eta$.
All witnesses centered in $Q_B$ are selected. Restricting targets cannot
lower the true infimum. Its measured defect is at least $\eta-14e$.
Reserve another $6e$ for numerical enclosures and use threshold $\eta/2$,
with $20e<\eta/4$. This accepts exactly the true periods. Strict cutoff
collars ensure that every generator in Lemma~\ref{lem:recognition}
is present. The same tests on pairs of sources give their primitive
orbit relation: a period translating the first center to the second
must translate the first body to the second body, since Steiner points
lie in the interiors and different components are disjoint. All
primitive orbits are represented, so their number is recovered.

Accepted true vectors have a prior bound $U_0$ on their Euclidean
lengths. A nonzero determinant of two such periods is an integer
multiple of $\covol\Pi\ge V_*$. Determinant perturbations are
$O(U_0e)$; hence independence is decided below this gap. Select an
independent true pair with matrix $D$. Its generated sublattice has
index
\[
       n=|\det D|/\covol\Pi\le Q:=\max(1,\lceil U_0^2/V_*\rceil).
\]
Every coordinate of $D^{-1}w_j$ has reduced denominator dividing $n$
and a prior-bounded numerator. The inverse-matrix identity bounds its
estimation error by $C_*e$. Distinct rationals of denominators at most
$Q$ are separated by at least $Q^{-2}$. Finite search locks the exact
coordinates when $C_*e<1/(3Q^2)$.

Let $q$ be their common denominator, which divides $n$. A column
Hermite basis $H_0$ of the integer span of $qI_2$ and $qD^{-1}w_j$
gives
\begin{equation}\label{eq:period-basis}
                         \Pi=(DH_0/q)\Z^2.
\end{equation}
The integer group contains $q\Z^2$, so its Hermite entries have a
prior bound. Replacing $D$ by its estimate gives a matched true basis
with error $O(e)$. We have recovered exact integer relations, not
exact real coordinates. The primitive free area is
\[
             A_* =\covol\Pi-\sum_{j=1}^{r_*}\area(C_j).
\]
Uniform perimeter bounds and the parallel-body inclusions give area
error $O(e)$ for each matched component, proving the last assertion.
Support suprema and center integrals used in the tests admit certified
finite angular meshes, because the bounded bodies have a uniform
support Lipschitz bound. Their numerical errors are charged above.
\end{proof}

\begin{proof}[Proof of Theorem~\ref{thm:adaptive}]
Choose a sufficiently large fixed target square so that every component
needed in Lemma~\ref{lem:locking} and every coarse bracket is protected.
Separation bounds their number. Choose the fixed observation and launch
apertures as in Lemma~\ref{lem:stopped}. Apply
Proposition~\ref{prop:patch} with an integer angular mesh satisfying
$h\asymp\nu^{1/(s-2)}$ and a sufficiently small constant factor.
Its reconstructed bodies have matched $C^2$ support error $O(\nu)$
and Hausdorff error $e=O(h^s)$. For small $\nu$, all thresholds in
Lemma~\ref{lem:locking} hold. It gives exact primitive discrete data
and period-basis error $O(h^s)$. Repeating one recovered body per
primitive orbit under this basis gives a physical periodic presentation.
Indeed the inverse basis remains bounded, so only a bounded number of
integer translates can enter a bounded fundamental patch. The positive
separation and curvature margins then persist. The reconstructed
bodies have area error $O(h^s)$ and the covolume has the same order;
hence the primitive free-area error is also $O(h^s)$, in particular
$O(\nu)$.

The query count of Proposition~\ref{prop:patch} and its fixed local
stencil size give \eqref{eq:centers}. Equation~\eqref{eq:patch-cost}
gives \eqref{eq:attempts}, since $s$ is fixed. The calibration bias
allocation in Proposition~\ref{prop:query} gives \eqref{eq:precision}.
Every preparation has already been charged there. All steps are finite,
and the output is a finite smooth representation with certified
numerical tolerances. This proves the theorem.
\end{proof}

Without a known $\eta$, successive refinements give the corresponding
pointwise eventual decisions for each fixed table, but not a uniformly
valid observable stopping certificate. The exact functional inverse
combined with Lemma~\ref{lem:recognition} needs no supplied positive
margin: it tests zeros exactly. This distinction, as well as the
bounded periodic presentation itself, is retained. No inference of
crystallinity from an arbitrary nonperiodic configuration is made.
